\documentclass[10pt,a4paper]{amsart}
\usepackage[margin=19mm]{geometry}
\usepackage{amsmath,amssymb,mathtools}
\usepackage[scr=rsfs,cal=euler]{mathalfa}
\usepackage{xfrac}
\usepackage[unicode,hidelinks,bookmarks=false]{hyperref}
\newcommand{\ie}{i.e.\ }
\newcommand{\cf}{cf.\ }
\newcommand{\resp}{respectively\ }
\newcommand{\Subsection}{\subsection}
\providecommand{\nobreakdash}{\nobreak\hbox{-}}
\setlength{\emergencystretch}{2em}
\multlinegap0pt
\usepackage{bm}
\usepackage{bbm}
\usepackage{dsfont}
\usepackage{xspace}
\usepackage{cancel}
\usepackage{mathtools}
\usepackage{amsthm}
\makeatletter
\@ifundefined{theorem}{\newtheorem{theorem}{Theorem}}{}
\@ifundefined{lemma}{\newtheorem{lemma}[theorem]{Lemma}}{}
\makeatother
\newcommand{\CH}{\mathsf{CH}}
\newcommand{\forceP}{\mathbb{P}}
\newcommand{\forceR}{\mathbb{R}}
\title[Forcing upper $\Sigma$-uniformization in the presence of low]{Forcing upper $\Sigma$-uniformization in the presence of lower $\Pi$-reduction or uniformization}
\author{Stefan Hoffelner}
\date{Source: arXiv:2511.05081v2 (CC BY 4.0)}
\begin{document}
\maketitle

\noindent\emph{Source and licence.} Forcing upper $\Sigma$-uniformization in the presence of lower $\Pi$-reduction or uniformization. Version arXiv:2511.05081v2 (CC BY 4.0), distributed under the Creative Commons Attribution 4.0 International licence, as recorded in the arXiv licence metadata for this exact version and shown on the abstract page https://arxiv.org/abs/2511.05081v2. Full rights notice: licenses/arxiv-2511.05081-CC-BY-4.0.txt.\par\smallskip

\noindent\textit{Editorial scope.} Counted unit: the Theorem whose source label is \texttt{None} in the article (source0.tex lines 2640--2642, source label \texttt{None}), delivered together with the article's own complete original proof of it (source0.tex lines 2644--2735, one \texttt{proof} environment of 92 lines). No mathematics is written, repaired or re-derived: statement and proof are the article's own text line for line. Required context retained uncounted: UniformizationClosureLemma (lines 2296--2307); UniformizationPersistenceTailLemma (lines 2313--2337); NoUnwantedCodesSecondMainTheorem (lines 2624--2630). Every definition, display and label of those lines is retained unchanged. Omitted from this package and not redistributed: 1--2295, 2308--2312, 2338--2623, 2631--2639, 2643--2643, 2736--2851. Nothing inside a retained excerpt is omitted or altered: every retained body is delivered exactly as the article's lines are written, so no omission receipt of any kind accompanies this package. Citations: the retained excerpts carry no citation command at all, so no bibliography entry of the article is transcribed and no reference is redistributed. Reference adaptation: none was needed and none was made; every reference inside the delivered unit resolves inside it, so no unresolved reference or citation is produced. Printed slips kept literal: no printed word, symbol, hypothesis or number of the counted statement or of its proof was altered. Layout and class: the article's own class is \texttt{article} with packages including inputenc, hyperref, boxedminipage, amsfonts, amsmath, amssymb, graphicx, amsthm, t1enc, subfig, cancel, textcmds, mathabx; the wrapper is the lane's pinned \texttt{amsart} editorial wrapper at 10pt on A4, which restates the theorem-like environments the retained text needs and adds the article's own macro definitions that the retained text needs (\texttt{\textbackslash CH}, \texttt{\textbackslash forceP}, \texttt{\textbackslash forceR}), with extra wrapper packages (bm, bbm, dsfont, xspace, cancel, mathtools). The delivered numbering is the unit's own. Assets: no retained span contains a figure environment, a graphics-inclusion command, a PDF-page inclusion, a table or a TikZ picture; no image file is copied into this package, the package declares no asset dependency and no image file is redistributed with it. Licence: version arXiv:2511.05081v2 of this article is distributed under the Creative Commons Attribution 4.0 International licence, as recorded in the arXiv licence metadata for this exact version and shown on the abstract page https://arxiv.org/abs/2511.05081v2. A full rights notice is delivered at licenses/arxiv-2511.05081-CC-BY-4.0.txt. Characters: every retained excerpt is pure ASCII, so the delivered engine, XeLaTeX with its default Latin Modern text font, reports no missing character.
\begin{lemma}\label{UniformizationClosureLemma}
Let $(\forceP,\dot I)\in\mathcal C_\infty$. Then:
\begin{enumerate}
\item $C^\omega(\forceP)$ belongs to the same current fixed-point class $\mathcal C_\infty$;
\item it avoids the same forbidden family;
\item its auxiliary tentative-value set is the union of the tentative-value tuples carried by $\forceP$ and by all witness forcings inserted during the closure, with all global ranks and stored keys unchanged;
\item if an actual factor of $C^\omega(\forceP)$ codes $(m,x,y)$ and, when that factor appears, some forcing in the current fixed-point class forces $(x,y)\notin A_m$, then
\[
C^\omega(\forceP)\Vdash(x,y)\notin A_m.
\]
\end{enumerate}
\end{lemma}

\begin{lemma}\label{UniformizationPersistenceTailLemma}
Suppose that the tuple
\[
(m,x,y,\rho,\kappa),\qquad \rho=(\xi,\zeta),
\]
is inserted into the auxiliary tentative-value set because
\[
(x,y)\in A_m
\]
and no forcing in the class against which persistence is tested at rank $\rho$ forces
\[
(x,y)\notin A_m.
\]
Then the following hold throughout the later main construction.
\begin{enumerate}
\item Every later fixed-point class is a subclass of the class against which the persistence of $y$ was tested.
\item Every later stage block and every bounded quotient of the main iteration belongs, after translating each admissible support together with its local names, to the corresponding earlier class. This remains true for products and amalgamations, for the closures $C^\omega(\forceP)$, and after later restarts.
\item No later bounded quotient forces $(x,y)\notin A_m$.
\item Consequently,
\[
(x,y)\in A_m
\]
in the final model.
\end{enumerate}
\end{lemma}

\begin{lemma}\label{NoUnwantedCodesSecondMainTheorem}
The no-unwanted-codes lemma applies to the allowable classes used in this section. In particular, in the final model a finite tuple $\vec z$ is coded into $\vec S$ if and only if an actual factor
\[
\operatorname{Code}(\vec z,\eta)
\]
occurs in the main iteration or in one of its allowable stage blocks.
\end{lemma}

\begin{theorem}
Let $G_{\omega_1}$ be a $\forceP_{\omega_1}$-generic filter over $W$. Then in $W[G_{\omega_1}]$ the $\Pi^1_3$-uniformization property holds and $\Sigma^1_n$-uniformization holds for every $n\geq4$. Moreover, the reals admit a good $\Sigma^1_5$-well-order.
\end{theorem}

\begin{proof}
For a tuple $(m,x,y)$ let
\[
\operatorname{Coded}_{\vec S}(m,x,y)
\]
abbreviate the $\Sigma^1_3$-statement from Section~3 saying that $(m,x,y)$ is coded into $\vec S$. If $m$ is the G\"odel number of a $\Pi^1_3$-set $A_m$ in the plane, define
\[
f_m(x)=y
\quad\Longleftrightarrow\quad
(x,y)\in A_m
\ \text{ and }\
\neg\operatorname{Coded}_{\vec S}(m,x,y).
\]
This is a $\Pi^1_3$ definition.

Fix $m$ and a real $x$ in the final model. We distinguish two cases.

\medskip
\noindent
\emph{A tentative value occurs.}
Assume that
\[
(m,x,y,\rho,\kappa)\in I_{\omega_1}.
\]
Choose the least pair $(\rho_0,\kappa_0)$ occurring for $(m,x)$ and let $y_0$ be the corresponding real, using the least supported presentation in case of a tie. We claim first that $(m,x,y_0)$ was not coded before $y_0$ became the selected value.

Suppose otherwise and consider the first actual factor which coded $(m,x,y_0)$. There are two possibilities. If it was produced directly by a persistent-value stage of the outer iteration, then the value selected at that stage had a smaller rank--key pair than $(\rho_0,\kappa_0)$; if the selected value was $y_0$ itself, the stage did not code $(m,x,y_0)$. The first alternative contradicts the minimality of $(\rho_0,\kappa_0)$ and the second contradicts the choice of the factor.

It remains to consider a factor occurring inside a closed forceable-out block. Let $N_0$ be the model over which that allowable block is computed and let $A$ be the admissible support of the local bookkeeping entry which produced the factor. Then $y_0\in N_0[G_A]$. If a forcing in the current fixed-point class could force $(x,y_0)\notin A_m$, Lemma~\ref{UniformizationClosureLemma} would make the closed block force
\[
(x,y_0)\notin A_m,
\]
contrary to Lemma~\ref{UniformizationPersistenceTailLemma}. If no such force-out forcing exists, then, at the fixed point, one further application of the support-local derivative selects a persistent value from the same candidate universe $N_0[G_A]$. It either selects $y_0$, in which case $(m,x,y_0)$ is left uncoded, or selects a tentative value with a smaller rank--key pair, which is carried into $I_{\omega_1}$. Both alternatives are impossible. Thus no factor codes $(m,x,y_0)$ before it is selected.

After $y_0$ has been selected, Lemma~\ref{UniformizationPersistenceTailLemma} gives
\[
(x,y_0)\in A_m
\]
in the final model. Every later fixed-point class is a subclass of the class against which its persistence was tested. Products, closures and restarts translate the admissible support of its local presentation and preserve the stored pair $(\rho_0,\kappa_0)$; they do not repeat the candidate search in a larger support extension. Since this is the least pair which occurs in $I_{\omega_1}$ for $(m,x)$, the rank/key convention continues to select $y_0$ whenever $(m,x)$ is considered. Hence no later factor codes $(m,x,y_0)$. By Lemma~\ref{NoUnwantedCodesSecondMainTheorem},
\[
\neg\operatorname{Coded}_{\vec S}(m,x,y_0).
\]
Thus $f_m(x)=y_0$.

Let $y'\neq y_0$ be any real in the final model. After names for $m,x,y'$ have appeared, the tagged bookkeeping presents them at a later uniformization stage. The tentative value with least rank--key pair remains $y_0$, with rank $\rho_0$ and key $\kappa_0$, so the stage uses
\[
\operatorname{Code}(m,x,y',\eta)
\]
for a free coding area. Hence $y_0$ is the unique real satisfying the definition of $f_m(x)$.

\medskip
\noindent
\emph{No tentative value occurs.}
Assume that $I_{\omega_1}$ contains no tuple of the form
\[
(m,x,y,\rho,\kappa).
\]
Let $y$ be any real in the final model and choose a later uniformization stage after names for $m,x,y$ have appeared. If $(x,y)\notin A_m$ already holds at that stage, it remains true in the final model, since this is a $\Sigma^1_3$-statement. Otherwise the current fixed-point class carries no persistent tentative value for $(m,x)$. The forceable-out clause therefore chooses a forcing $\forceR_{m,x,y}$ with
\[
\forceR_{m,x,y}\Vdash(x,y)\notin A_m
\]
and uses
\[
C^\omega(\forceR_{m,x,y})
\]
as the stage forcing. Hence $(x,y)\notin A_m$ holds from that stage on. Since every real $y$ is eventually treated,
\[
(A_m)_x=\emptyset
\]
in the final model.

We have proved that $f_m$ has exactly one value on every nonempty section of $A_m$ and no value on an empty section. Since $m$ was arbitrary, $\Pi^1_3$-uniformization holds.

We next verify the well-order conclusion. Let
\[
<\;:=\bigcup_{\beta<\omega_1}<_\beta.
\]
Every successor step is an end-extension and every proper initial segment is countable. Every real in the final model is the interpretation of a name from some bounded outer stage, and the tagged bookkeeping presents that name later at a well-order stage. Hence every real enters the field of $<$. Since the final model satisfies $\CH$, the order has type $\omega_1$.

If $z$ codes the set of $<$-predecessors of $x$, then after $x$ enters the order this predecessor set never changes. The exact tuple $(1,1,x,z)$ is never put into a forbidden family, and the bookkeeping codes it cofinally often. Conversely, if $z$ does not code the correct predecessor set, choose a later well-order stage at which a supported name for $z$ is considered after $x$ has entered the order. The tuple $(1,1,x,z)$ is then added to the new forbidden family. It can be coded only below a bounded stage and is therefore not coded cofinally often. Consequently
\[
z\text{ codes the set of $<$-predecessors of }x
\quad\Longleftrightarrow\quad
(1,1,x,z)\text{ is coded cofinally often}.
\]
The right-hand side is a $\Sigma^1_5$-formula, as in the proof of the first main theorem. Thus $<$ is a good $\Sigma^1_5$-well-order.

Finally, $\Pi^1_3$-uniformization implies $\Sigma^1_4$-uniformization by uniformizing a $\Pi^1_3$ relation on pairs and projecting the selected pair. The good $\Sigma^1_5$-well-order gives $\Sigma^1_n$-uniformization for every $n\geq5$ by Addison's argument. Hence
\[
\Sigma^1_n\text{-uniformization holds for every }n\geq4.
\]
\end{proof}

\end{document}
